\expandafter\gdef\csname SNside1\endcsname{\node[anchor=north west,inner sep=0pt,text width=48mm,font=\SideFont,align=left] at ([xshift=138mm,yshift=-547.22pt]current page.north west) {\textcolor{SNink}{\textbf{课程连接 · 网络拓扑只是图的一种模型}}\par\vspace{1.5mm}计算机网络可把路由器视为顶点、链路视为边，但边权可能表示时延、费用或带宽的变换。先明确权的含义再选择算法；同一个物理网络可以对应不同优化问题。};}
\expandafter\gdef\csname SNside2\endcsname{\node[anchor=north west,inner sep=0pt,text width=48mm,font=\SideFont,align=left] at ([xshift=138mm,yshift=-333.03pt]current page.north west) {\textcolor{SNink}{\textbf{最小反例 · n减1条边不够判树}}\par\vspace{1.5mm}三点成环再加一个孤立点，共4顶点3边，恰有\(n-1\)条边却不连通且有环。判树还要加连通或无环条件，不能只由数量推出结构。\par\smallskip\begin{tikzpicture}[x=1mm,y=1mm,>=Stealth,node distance=5mm]
\node[draw=SNink!50,fill=SNblue,rounded corners=1mm,text width=34mm,align=center,inner sep=1.2mm,font=\fontsize{7.5}{9}\selectfont] (m0) at (0,0) {A};
\node[draw=SNink!50,fill=SNblue,rounded corners=1mm,text width=34mm,align=center,inner sep=1.2mm,font=\fontsize{7.5}{9}\selectfont] (m1) at (0,-11) {B};
\node[draw=SNink!50,fill=SNblue,rounded corners=1mm,text width=34mm,align=center,inner sep=1.2mm,font=\fontsize{7.5}{9}\selectfont] (m2) at (0,-22) {C};
\node[draw=SNink!50,fill=SNblue,rounded corners=1mm,text width=34mm,align=center,inner sep=1.2mm,font=\fontsize{7.5}{9}\selectfont] (m3) at (0,-33) {孤点D};
\draw[->,SNink] (m0) -- (m1);
\draw[->,SNink] (m1) -- (m2);
\draw[->,SNink] (m2) edge[bend left=40] (m0);
\end{tikzpicture}
\par{\fontsize{7.5}{9}\selectfont 4顶点3边：三角形加孤点，满足m=n−1却不是树。}};}
\expandafter\gdef\csname SNside3\endcsname{\node[anchor=north west,inner sep=0pt,text width=48mm,font=\SideFont,align=left] at ([xshift=138mm,yshift=-95.00pt]current page.north west) {\textcolor{SNink}{\textbf{错项排除 · 零权边不是没有边}}\par\vspace{1.5mm}带权邻接矩阵用无穷表示无边时，零是合法边权。若仍以if(A{[}u{]}{[}v{]})判断有边，就会漏掉零权边；AOV的0/1邻接矩阵与这种表示的语义不同。};}
\expandafter\gdef\csname SNside4\endcsname{\node[anchor=north west,inner sep=0pt,text width=48mm,font=\SideFont,align=left] at ([xshift=138mm,yshift=-95.00pt]current page.north west) {\textcolor{SNink}{\textbf{不变量 · 一条边结点参与两条链}}\par\vspace{1.5mm}邻接多重表让同一个无向边结点同时属于两个端点链。沿当前端点前进时选择其对应链接；这解释了为何可统一标记边，也解释了不能始终沿ilink。};}
\expandafter\gdef\csname SNside5\endcsname{\node[anchor=north west,inner sep=0pt,text width=48mm,font=\SideFont,align=left] at ([xshift=138mm,yshift=-95.00pt]current page.north west) {\textcolor{SNink}{\textbf{不变量 · DFS何时才允许回退}}\par\vspace{1.5mm}当前顶点仍有未访问邻点时必须深入，全部处理完才回退。输出相邻两点可以不直接相连，因为中间隐含了回退；仅检查输出相邻性会误排除合法DFS。};}
\expandafter\gdef\csname SNside6\endcsname{\node[anchor=north west,inner sep=0pt,text width=48mm,font=\SideFont,align=left] at ([xshift=138mm,yshift=-95.00pt]current page.north west) {\textcolor{SNink}{\textbf{最小反例 · 边最少而权更大}}\par\vspace{1.5mm}三点A、B、C，AB与BC权1，AC权3。BFS从A一跳发现C，少一条边却总权3；经B两跳总权2。因此BFS的分层证明不能覆盖不等权路径目标。};}
\expandafter\gdef\csname SNside7\endcsname{\node[anchor=north west,inner sep=0pt,text width=48mm,font=\SideFont,align=left] at ([xshift=138mm,yshift=-95.00pt]current page.north west) {\textcolor{SNink}{\textbf{错项排除 · 同权不等于MST必不唯一}}\par\vspace{1.5mm}若所有边权互异，MST一定唯一；逆命题不成立。一个本来就是树的连通图，即使边权相同，其生成树仍只有它本身，不能根据同权直接判不唯一。};}
\expandafter\gdef\csname SNside8\endcsname{\node[anchor=north west,inner sep=0pt,text width=48mm,font=\SideFont,align=left] at ([xshift=138mm,yshift=-95.00pt]current page.north west) {\textcolor{SNink}{\textbf{不变量 · Prim的key不是距离}}\par\vspace{1.5mm}Prim只问一条跨割边接新点的代价，key更新取边权；Dijkstra问源点到新点的完整路径，需加已确定前缀。把两者公式混写可能算出另一种对象。};}
\expandafter\gdef\csname SNside9\endcsname{\node[anchor=north west,inner sep=0pt,text width=48mm,font=\SideFont,align=left] at ([xshift=138mm,yshift=-95.00pt]current page.north west) {\textcolor{SNink}{\textbf{不变量 · 已确定与暂定要分开}}\par\vspace{1.5mm}Dijkstra中暂定dist只是一条已知路径的最好上界。选出最小暂定点后才永久确定；过程题若提前使用未来松弛的值，就把最终答案错误放入了中间状态。};}
\expandafter\gdef\csname SNside10\endcsname{\node[anchor=north west,inner sep=0pt,text width=48mm,font=\SideFont,align=left] at ([xshift=138mm,yshift=-95.00pt]current page.north west) {\textcolor{SNink}{\textbf{课程连接 · 同步前驱约束就是偏序}}\par\vspace{1.5mm}操作系统中多个事件的前驱关系可用DAG表示，零入度事件代表当前所有前驱已完成。拓扑排序只给合法先后，不计算耗时；带持续时间的工期还需AOE时间传播。};}
\expandafter\gdef\csname SNside11\endcsname{\node[anchor=north west,inner sep=0pt,text width=48mm,font=\SideFont,align=left] at ([xshift=138mm,yshift=-95.00pt]current page.north west) {\textcolor{SNink}{\textbf{错项排除 · 缩短一条关键活动未必够}}\par\vspace{1.5mm}存在两条等长关键路径时，只缩短其中一条独有活动，另一条仍保持原工期。要检查所有当前关键路径是否都被影响，并在缩短后重新计算新关键路径。};}
\expandafter\gdef\csname SNside12\endcsname{\node[anchor=north west,inner sep=0pt,text width=48mm,font=\SideFont,align=left] at ([xshift=138mm,yshift=-95.00pt]current page.north west) {\textcolor{SNink}{\textbf{错项排除 · 有负边与一定算错不同}}\par\vspace{1.5mm}负边意味着标准Dijkstra失去通用正确性保证，不代表每个含负边实例都会得错值。选项若用``只要有一条负边就必算错''的绝对语气，可用简单单路径反驳。};}
\expandafter\gdef\csname SNside13\endcsname{\node[anchor=north west,inner sep=0pt,text width=48mm,font=\SideFont,align=left] at ([xshift=138mm,yshift=-95.00pt]current page.north west) {\textcolor{SNink}{\textbf{错项排除 · 缩短一条关键活动未必够}}\par\vspace{1.5mm}存在两条等长关键路径时，只缩短其中一条独有活动，另一条仍保持原工期。要检查所有当前关键路径是否都被影响，并在缩短后重新计算新关键路径。};}
\expandafter\gdef\csname SNside14\endcsname{\node[anchor=north west,inner sep=0pt,text width=48mm,font=\SideFont,align=left] at ([xshift=138mm,yshift=-95.00pt]current page.north west) {\textcolor{SNink}{\textbf{课程连接 · 同步前驱约束就是偏序}}\par\vspace{1.5mm}操作系统中多个事件的前驱关系可用DAG表示，零入度事件代表当前所有前驱已完成。拓扑排序只给合法先后，不计算耗时；带持续时间的工期还需AOE时间传播。};}
\AddToShipoutPictureFG{\begin{tikzpicture}[remember picture,overlay]\ifcsname SNside\thepage\endcsname\csname SNside\thepage\endcsname\fi\end{tikzpicture}}